% Generated draft: requires mathematical and source review.
\documentclass[11pt]{article}
\usepackage[margin=25mm]{geometry}
\usepackage{fontspec}
\setmainfont{Latin Modern Roman}
\usepackage{amsmath,amssymb,mathtools,unicode-math}
\setmathfont{Latin Modern Math}
\usepackage{enumitem,needspace,xurl,hyperref}
\setlength{\parindent}{0pt}
\setlength{\parskip}{5pt}
\setlength{\emergencystretch}{3em}
\newcommand{\N}{\mathbb{N}}
\newcommand{\Z}{\mathbb{Z}}
\newcommand{\Q}{\mathbb{Q}}
\newcommand{\R}{\mathbb{R}}
\newcommand{\C}{\mathbb{C}}
\newcommand{\F}{\mathbb{F}}
\newcommand{\A}{\mathbb{A}}
\newcommand{\PP}{\mathbb{P}}
\newcommand{\E}{\mathbb{E}}
\newcommand{\eps}{\varepsilon}
\newcommand{\dd}{\,\mathrm{d}}
\begin{document}
\section*{UnsolvedMath ID: 10400119}
\textbf{Source problem code:} catalogue entry
\subsection{Definitions and mathematical statement}
Let \(\tau_\lambda(M)\) denote whichever quantum invariant of a three-manifold \(M\) is intended in the surrounding source, with \(\lambda\) representing its theory, level and any other parameters. The requested task is to find a three-dimensional topological interpretation of these values. Such an interpretation must identify topological data and a construction \(F_\lambda\) for which
\[F_\lambda(M)=\tau_\lambda(M)\]
for every admissible \(M\) and \(\lambda\), and explain intrinsically in three-dimensional topological terms what the construction computes. The displayed notation only expresses the required agreement; it does not supply the interpretation. Any orientation, framing, boundary or auxiliary structure required by the chosen quantum invariant must also be included and its role specified.

The research summary distinguishes Witten's Chern--Simons description, which it says provides a three-dimensional interpretation, from the \emph{purely topological} interpretation sought here. Thus the intended task is stronger than asserting that quantum invariants are unchanged by the relevant homeomorphisms. However, the excerpt gives no criterion for purely topological interpretation, no specified quantum-invariant family or target ring, and no criterion deciding whether a surgery, state-sum or field-theoretic construction meets that request. No such criterion or particular quantum theory is chosen in this draft.
\subsection{Short English statement}
Explain quantum three-manifold invariants through a construction with a genuinely three-dimensional topological meaning. The record contrasts this goal with the supplied Chern--Simons interpretation, but leaves the desired interpretation and invariant family undefined.
\subsection{Sources}
\begin{itemize}
\item Local UnsolvedMath registry, permanent ID \(10400119\), exact source code \texttt{catalogue entry}.
\item Ohtsuki, \emph{Problems on invariants of knots and 3-manifolds} (2002), Problem 7.4, PDF page 101: \url{https://msp.org/gtm/2002/04/gtm-2002-04-024s.pdf}. Registry access date 2026-07-29, pdf-font-structure extraction.
\item The research summary mentions Witten's Chern--Simons theory without a specific bibliographic reference.
\end{itemize}
\end{document}
